% Generated by theory/supplement/render_latex.py -- do not edit by hand.
% The prose and every formula live in theory/supplement/theorem1_document.py,
% which also drives the Word build, so the two outputs cannot diverge.
\documentclass[11pt,a4paper]{article}

\usepackage[T1]{fontenc}
\usepackage[utf8]{inputenc}
\usepackage[margin=1in]{geometry}
\usepackage{amsmath,amssymb,amsthm}
\usepackage{booktabs}
\usepackage{longtable}
\usepackage{array}
\usepackage{ragged2e}
\usepackage[table]{xcolor}
\usepackage{fancyhdr}
\usepackage{titlesec}
\usepackage{enumitem}
\usepackage[most]{tcolorbox}
\usepackage[hidelinks]{hyperref}

\definecolor{dbnavy}{HTML}{17324D}
\definecolor{dbblue}{HTML}{2A6F97}
\definecolor{dbmuted}{HTML}{5C6B78}
\definecolor{dbpale}{HTML}{EAF2F8}
\definecolor{dbcalloutbg}{HTML}{F4F8FB}

\titleformat{\section}{\normalfont\Large\bfseries\color{dbnavy}}{\thesection}{0.6em}{}
\titleformat{\subsection}{\normalfont\large\bfseries\color{dbblue}}{\thesubsection}{0.6em}{}
\titleformat{\subsubsection}{\normalfont\normalsize\bfseries\color{dbnavy}}{}{0em}{}

\pagestyle{fancy}
\fancyhf{}
\fancyhead[L]{\footnotesize\color{dbmuted}DBACT THEOREM 1  /  PROOF SUPPLEMENT}
\fancyhead[R]{\footnotesize\color{dbmuted}LOCAL - CONDITIONAL - AUDITABLE}
\fancyfoot[L]{\footnotesize\color{dbmuted}Baseline HEAD 64e1905e092d}
\fancyfoot[R]{\footnotesize\color{dbmuted}\thepage}
\renewcommand{\headrulewidth}{0.4pt}

\newtcolorbox{dbcallout}[1]{%
  breakable, enhanced, colback=dbcalloutbg, colframe=dbcalloutbg,
  boxrule=0pt, leftrule=2.5pt, sharp corners,
  borderline west={2.5pt}{0pt}{dbblue},
  left=9pt, right=9pt, top=7pt, bottom=7pt,
  title={#1}, coltitle=dbblue, colbacktitle=dbcalloutbg,
  fonttitle=\footnotesize\bfseries, titlerule=0pt, attach title to upper=\par\smallskip
}

\setlength{\parindent}{0pt}
\setlength{\parskip}{0.55em}
% A little slack for the occasional line holding a long symbol run or a DOI,
% rather than letting it stick out into the margin.
\setlength{\emergencystretch}{3em}

\begin{document}

\begin{titlepage}
\centering
\vspace*{3cm}
{\color{dbblue}\large\bfseries DBACT THEOREM 1}\par\medskip
{\color{dbnavy}\Huge\bfseries Proofs and Supporting Results}\par\bigskip
{\color{dbmuted}\large Boundary-measure-induced Local-CVT practical stability near a selected nondegenerate local branch}\par
\vspace{1.5cm}
\rowcolors{1}{dbpale}{dbpale}
\begin{tabular}{>{\bfseries\color{dbnavy}}p{3cm} >{\RaggedRight\arraybackslash}p{10cm}}
Branch & DBACT-research-v3 \\
Baseline HEAD & 64e1905e092d19fd9dcbf0e77fbd39f55f0562ba \\
Proof status & Conditional local theorem; A8 interval certificate not supplied \\
Wolfram & Connector exact checks passed; 9 local scripts exit 0 \\
Regression & 477 passed, 3 skipped \\
\end{tabular}
\vspace{1.5cm}

{\color{dbblue}\footnotesize\bfseries CLAIM BOUNDARY}\par\medskip
{\color{dbnavy}\bfseries No global PL. No formal caging. No finite-time phase completion. No unconditional safety or cooperative-transport success.}
\end{titlepage}

\tableofcontents
\newpage

\section{Scope and frozen claim}

Bounded local-map disagreement and bounded reference variation imply practical stability of the boundary-induced Local-CVT allocation dynamics around a nondegenerate local CVT branch.

The agents are first-order, $\dot p_i = u_i$, and the nominal ENCLOSE law is $u_i^{\mathrm{nom}} = -k_c\,(p_i - \widehat c_i)$. The ideal objective is the \emph{saturated} limited-range coverage cost

\[
F_R(r) = \min\{r^2,\ R_\ell^2\},\qquad H^*(P,t) = \int_D \min_i F_R\bigl(\|q-p_i\|\bigr)\,\phi^*(q,t)\,dq
\]

Saturation is not cosmetic. With the unsaturated truncated cost $\int_{V_i^\ell}\|q-p_i\|^2\phi^*\,dq$ the integration domain moves with $p_i$, and the Reynolds boundary term on the circle $\|q-p_i\| = R_\ell$ does not cancel; the standard CVT gradient identity is then false. Section 5 shows the cancellation that saturation buys.

\subsection{Deliberate exclusions}

This supplement does not prove any of the following.

\begin{itemize}[leftmargin=1.2em,itemsep=2pt]
  \item global CVT optimality, or convergence from arbitrary initial configurations;
  \item finite-time SEARCH-to-HOLD completion;
  \item geometric or topological caging;
  \item contact mechanics, or unconditional cooperative-transport success;
  \item CBF forward invariance when feasibility, speed-cap, or deviation contracts are violated;
  \item a uniform reduced-Hessian bound for an unspecified local branch.
\end{itemize}

Collision safety is a separate result and is deliberately not a dependency of Theorem 1.

\newpage

\section{Notation}

$D \subset \mathbb{R}^2$ is a fixed compact rectangle of width $W_D$ and height $H_D$; $P = (p_1,\dots,p_N) \in \mathbb{R}^{2N}$ stacks the agent positions. All spatial norms are Euclidean, and the stacked norm is the Euclidean product norm

\[
\|d\|^2 = \sum_{i=1}^{N}\|d_i\|^2
\]

$\Gamma(t)$ is the true object boundary, with perimeter $P_\Gamma \le P_{\max}$, and $\xi(s,t)$ is the offset-boundary target parameterised by arc length $s$. The ideal boundary measure is the pushforward of arc length, and the ideal density is its Gaussian convolution above a positive floor:

\[
\nu_t^* = \xi(\cdot,t)_{\#}\,ds,\qquad \phi^*(q,t) = \phi_0 + \bigl(K_\sigma * \nu_t^*\bigr)(q) = \phi_0 + \int_{\Gamma(t)} K_\sigma\bigl(q-\xi(s,t)\bigr)\,ds
\]

The kernel is the \emph{unnormalised} Gaussian $K_\sigma(x) = \exp\bigl(-\|x\|^2/(2\sigma^2)\bigr)$. The local cell, its mass and its centroid are

\[
V_i^\ell(P) = \mathrm{Vor}_i(P)\cap D\cap B(p_i,R_\ell),\qquad m_i^* = \int_{V_i^\ell}\phi^*\,dq,\qquad c_i^* = \frac{1}{m_i^*}\int_{V_i^\ell} q\,\phi^*(q,t)\,dq
\]

$\widehat\phi_i$ is agent $i$'s mapped density, $\widehat c_i$ its numerically computed centroid, $\mathcal{C}^*(t)$ the selected local minimiser branch inside a fixed compact regular tube $X$, and

\[
V(P,t) = H^*(P,t) - H^*_{\min}(t) \ \ge\ 0,\qquad H^*_{\min}(t) = \min_{P\in X} H^*(P,t)
\]

\subsection{Kernel constants}

These three constants are used throughout; all are Wolfram-verified (entries W-RR-01 to W-RR-03) and derived in Step 1 of Lemma 1.

\begin{equation*}
\|K_\sigma\|_1 = 2\pi\sigma^2,\qquad \|\nabla K_\sigma\|_1 = \sqrt{2}\,\pi^{3/2}\sigma,\qquad \|\nabla K_\sigma\|_\infty = \frac{e^{-1/2}}{\sigma}
\tag{K}
\end{equation*}

\subsection{Units}

Arc length is in metres and $K_\sigma$ is dimensionless, so $\phi$ has units of m, an $L^1$ density error and a cell mass have units $\mathrm{m}^3$, a shifted first moment $\mathrm{m}^4$, $H^*$ has $\mathrm{m}^5$, $\nabla H^*$ has $\mathrm{m}^4$, and a command disturbance has $\mathrm{m}/\mathrm{s}$. The ledger in Appendix A carries the unit of every constant so that geometric length and density mass are never silently added.

\newpage

\section{Assumptions A1-A10}

\begin{center}
\small
\begin{longtable}{>{\bfseries}p{0.9cm} >{\RaggedRight\arraybackslash}p{8.4cm} >{\RaggedRight\arraybackslash}p{4.6cm}}
\toprule
\textbf{ID} & \textbf{Assumption} & \textbf{How it is discharged} \\
\midrule
\endhead
A1 & $D$ is a fixed rectangle and the solution remains in a fixed compact configuration tube $X$ on the claimed interval. & assumption / invariant-domain check \\
A2 & Pairwise separation is at least $d_s>0$, and $R_{\mathrm{comm}} \ge 2R_\ell$. & runtime certificate + config validator (fail-closed) \\
A3 & $\Gamma(t)$ is closed and piecewise $C^2$ with $P_\Gamma \le P_{\max}$; version 1 uses a rigid boundary with bounded twist. & assumption / scenario certificate \\
A4 & $\phi_0>0$, $\sigma>0$, and the unnormalised Gaussian above is used. & analytic + config \\
A5 & Map weights are nonnegative and a certificate supplies $\varepsilon_b,\varepsilon_n,\varepsilon_d$, voxel $v$, $h_s$, $E_w$, $M_{\mathrm{drop}}$ and $s_\sigma$ as upper bounds. & runtime / offline certificate \\
A6 & Local CVT uses a rectangular midpoint grid of side at most $h$, and $\eta_m < m_{\min}$. & code contract + fail-closed test \\
A7 & Gap, frontier and transport density biases are disabled in theorem mode, or their stacked command effect is included in $\delta$. & config + perturbation ledger \\
A8 & On a symmetry-reduced convex chart of the selected branch, $\nabla^2 H^*_{\mathrm{red}}(z,t) \succeq m_H I$ uniformly, $m_H>0$. & CONDITIONAL - no interval certificate supplied \\
A9 & Voronoi ties and the threshold set $\|q-p_i\|=R_\ell$ have measure zero, and the solution stays on one regular branch tube. & local theorem assumption \\
A10 & ZOH, numerical, tracking and safety-filter deviations from the ideal field have certified stacked-norm bounds. & analytic + runtime certificate \\
\bottomrule
\end{longtable}
\end{center}

\begin{dbcallout}{STATUS OF THE ASSUMPTION SET}
The implication A1-A10 $\Rightarrow$ Theorem 1 is proved. A5, A8, A9 and A10 are not facts about every run: they are explicit assumptions and certificate gates. The result is therefore a complete conditional local theorem, not an unconditional system-success claim.
\end{dbcallout}

\newpage

\section{Lemma 1 - boundary and map to density consistency}

\subsection{Statement}

Assume A3-A5, let the arc partition satisfy $0 \le \Delta s_k \le h_s$, and suppose every source discarded by the restriction step is at least $s_\sigma\sigma$ away from every query point in $B(p_i,R_\ell)\cap D$. Define the target, weight and tail errors

\[
\varepsilon_{\xi,i} = \varepsilon_{b,i} + \varepsilon_{d,i} + 2 d_{\max}\sin\!\left(\frac{\varepsilon_{n,i}}{2}\right),\qquad \varepsilon_{x,i} = \varepsilon_{\xi,i} + \sqrt{2}\,v
\]

\[
E_{w,i} = L_{\mathrm{miss},i} + M_{\mathrm{spur},i} + \sum_{k\,\in\,\mathrm{matched}} \bigl|\widehat w_{ik} - \Delta s_k\bigr|
\]

Then

\begin{equation*}
\bigl\|\widehat\phi_i - \phi^*\bigr\|_{L^1(B(p_i,R_\ell)\cap D)} \ \le\ \varepsilon_{\phi,i} \ :=\ L_{K,1} P_{\max}\!\left(\varepsilon_{\xi,i} + \sqrt{2}\,v + \frac{L_T h_s}{2}\right) + 2\pi\sigma^2 E_{w,i} + \varepsilon_{\mathrm{tail},i}
\tag{L1}
\end{equation*}

with $L_{K,1} := \|\nabla K_\sigma\|_1$ and

\[
\varepsilon_{\mathrm{tail},i} \ \le\ \pi R_\ell^2\,M_{\mathrm{drop},i}\,e^{-s_\sigma^2/2}
\]

\begin{dbcallout}{WHY THREE MEASURES}
The comparison is made after convolution, through an exact discrete intermediate measure. It is never claimed that a continuous arc-length measure converges to a sum of Dirac masses in total variation - that statement is false.
\end{dbcallout}

\[
\nu_t^*\ \longrightarrow\ \nu_{h,t}^* = \sum_k \Delta s_k\,\delta_{\xi_k^*}\ \longrightarrow\ \widehat\nu_{i,t} = \sum_k \widehat w_{ik}\,\delta_{\widehat\xi_{ik}}
\]

\subsection{Proof}

\subsubsection*{Step 1. The three kernel constants}

The mass separates into one-dimensional Gaussians:

\[
\|K_\sigma\|_1 = \int_{\mathbb{R}^2} e^{-\|x\|^2/(2\sigma^2)}dx = \left(\int_{\mathbb{R}} e^{-x^2/(2\sigma^2)}dx\right)^{\!2} = \bigl(\sigma\sqrt{2\pi}\bigr)^2 = 2\pi\sigma^2
\]

Since $\nabla K_\sigma(x) = -\,\sigma^{-2} x\,K_\sigma(x)$, its magnitude is radial, $\|\nabla K_\sigma(x)\| = (r/\sigma^2)e^{-r^2/(2\sigma^2)}$ with $r = \|x\|$. Integrating in polar coordinates,

\[
\|\nabla K_\sigma\|_1 = \int_0^{\infty}\frac{r}{\sigma^2}e^{-r^2/(2\sigma^2)}\,2\pi r\,dr = \frac{2\pi}{\sigma^2}\int_0^{\infty} r^2 e^{-r^2/(2\sigma^2)}dr = \frac{2\pi}{\sigma^2}\cdot\sigma^3\sqrt{\frac{\pi}{2}} = \sqrt{2}\,\pi^{3/2}\sigma
\]

For the sup-norm, write $g(r) = (r/\sigma^2)e^{-r^2/(2\sigma^2)}$. Then

\[
g'(r) = \frac{1}{\sigma^2}e^{-r^2/(2\sigma^2)}\left(1 - \frac{r^2}{\sigma^2}\right) = 0 \ \Longleftrightarrow\ r = \sigma \quad (r>0)
\]

and $g(0) = 0$, $g(r)\to 0$ as $r\to\infty$, so the interior stationary point is the maximum:

\[
\|\nabla K_\sigma\|_\infty = g(\sigma) = \frac{\sigma}{\sigma^2}e^{-1/2} = \frac{e^{-1/2}}{\sigma}
\]

\subsubsection*{Step 2. Translation inequality in $L^1$}

This is the workhorse of the whole lemma: moving a source by a distance costs at most $\|\nabla K_\sigma\|_1$ times that distance, in $L^1$. For $a,b \in \mathbb{R}^2$ the fundamental theorem of calculus along the segment $\tau \mapsto b + \tau(a-b)$ gives

\[
K_\sigma(x-a) - K_\sigma(x-b) = -\int_0^1 \nabla K_\sigma\bigl(x - b - \tau(a-b)\bigr)\cdot(a-b)\,d\tau
\]

Taking absolute values, integrating in $x$ and exchanging the order of integration by Tonelli (the integrand is nonnegative),

\begin{equation*}
\bigl\|K_\sigma(\cdot-a)-K_\sigma(\cdot-b)\bigr\|_1 \le \|a-b\|\int_0^1\!\!\int_{\mathbb{R}^2}\bigl\|\nabla K_\sigma\bigl(x-b-\tau(a-b)\bigr)\bigr\|\,dx\,d\tau = \|\nabla K_\sigma\|_1\,\|a-b\|
\tag{T}
\end{equation*}

the inner integral being translation-invariant and therefore equal to $\|\nabla K_\sigma\|_1$ for every $\tau$.

\subsubsection*{Step 3. Target error}

With $\xi = b + d\,n$ and the observed $\widehat\xi = \widehat b + \widehat d\,\widehat n$, the triangle inequality gives

\[
\|\widehat\xi - \xi\| \le \|\widehat b - b\| + |\widehat d - d| + d_{\max}\,\|\widehat n - n\|
\]

Both normals are unit vectors, so if they subtend an angle $\theta$ then

\[
\|\widehat n - n\| = \sqrt{2 - 2\cos\theta} = 2\sin\frac{\theta}{2},\qquad 0\le\theta\le\pi
\]

which is monotone in $\theta$; substituting $\theta \le \varepsilon_n$ yields $\varepsilon_\xi$. Replacing a target by a representative of the same square voxel of side $v$ moves it by at most the voxel diagonal $\sqrt2\,v$, which gives $\varepsilon_x = \varepsilon_\xi + \sqrt2\,v$.

\begin{dbcallout}{SCOPE OF THE VOXEL TERM}
$\sqrt2\,v$ is valid only while the fused representative stays inside its own voxel. If the fusion law can move it outside, the bound must be re-derived from that law; it must not be reused as written.
\end{dbcallout}

\subsubsection*{Step 4. Curve quadrature}

Compare the exact discrete measure with the continuous one. On arc $k$, of length $\Delta s_k$ and with representative $\xi_k^*$,

\[
K_\sigma * \nu_h^* - K_\sigma * \nu^* = \sum_k \int_{\mathrm{arc}_k}\Bigl[K_\sigma(\cdot-\xi_k^*) - K_\sigma\bigl(\cdot-\xi(s)\bigr)\Bigr]ds
\]

Apply (T) inside the sum. If $\xi$ is $L_T$-Lipschitz in arc length then $\|\xi_k^* - \xi(s)\| \le L_T|s - s_k|$, and for any representative in the arc $\int_{\mathrm{arc}_k}|s-s_k|\,ds \le \Delta s_k^2/2$. Hence

\[
\bigl\|K_\sigma * \nu_h^* - K_\sigma * \nu^*\bigr\|_1 \le L_{K,1}L_T\sum_k \frac{\Delta s_k^2}{2}
\]

Because $0 \le \Delta s_k \le h_s$ we have $\Delta s_k^2 \le h_s\Delta s_k$, so the sum telescopes into the perimeter (Wolfram entry W-RR-05):

\[
\sum_k \Delta s_k^2 \ \le\ h_s\sum_k \Delta s_k \ =\ h_s P_\Gamma \ \le\ h_s P_{\max} \quad\Longrightarrow\quad \bigl\|K_\sigma * \nu_h^* - K_\sigma * \nu^*\bigr\|_1 \le \tfrac12 L_{K,1}L_T P_{\max} h_s
\]

Choosing the arc midpoint would give the sharper factor $\Delta s_k^2/4$; the certificate uses the conservative $1/2$ so that it stays valid for any registered representative inside the arc.

\subsubsection*{Step 5. Map displacement and weight mismatch}

Match the observed atoms to the exact ones and split the difference into a displacement part, a matched-weight part, and missing/spurious mass:

\[
\widehat\nu_i - \nu_h^* = \underbrace{\sum_{\mathrm{matched}}\Delta s_k\bigl(\delta_{\widehat\xi_k}-\delta_{\xi_k^*}\bigr)}_{\text{displacement}} + \underbrace{\sum_{\mathrm{matched}}\bigl(\widehat w_k-\Delta s_k\bigr)\delta_{\widehat\xi_k}}_{\text{weight}} - \sum_{\mathrm{missed}}\Delta s_k\,\delta_{\xi_k^*} + \sum_{\mathrm{spurious}}\widehat w_k\,\delta_{\widehat\xi_k}
\]

Convolve and bound each group. By (T), the displacement group costs

\[
\Bigl\|\sum_{\mathrm{matched}}\Delta s_k\bigl(K_\sigma(\cdot-\widehat\xi_k)-K_\sigma(\cdot-\xi_k^*)\bigr)\Bigr\|_1 \le L_{K,1}\,\varepsilon_x \sum_k \Delta s_k \le L_{K,1}P_{\max}\,\varepsilon_x
\]

The remaining three groups are bounded by their total variation times the kernel mass, because $\|w\,K_\sigma(\cdot-y)\|_1 = |w|\,\|K_\sigma\|_1$:

\[
\bigl\|K_\sigma * (\text{weight} + \text{missed} + \text{spurious})\bigr\|_1 \le \|K_\sigma\|_1\Bigl(L_{\mathrm{miss}} + M_{\mathrm{spur}} + \sum_{\mathrm{matched}}|\widehat w_k - \Delta s_k|\Bigr) = 2\pi\sigma^2 E_w
\]

\subsubsection*{Step 6. Restriction tail}

The restriction step discards far sources. It does not do so for free: the Gaussian has no compact support. A discarded atom of weight $\widehat w_k$ lies at distance at least $s_\sigma\sigma$ from every query $q$ in the disk, so pointwise

\[
K_\sigma(q-\widehat\xi_k) = \exp\!\left(-\frac{\|q-\widehat\xi_k\|^2}{2\sigma^2}\right) \le \exp\!\left(-\frac{s_\sigma^2}{2}\right)
\]

Summing the discarded weights and integrating over an area no larger than $\pi R_\ell^2$,

\[
\varepsilon_{\mathrm{tail},i} = \int_{B(p_i,R_\ell)\cap D}\sum_{\mathrm{dropped}}\widehat w_k K_\sigma(q-\widehat\xi_k)\,dq \le \pi R_\ell^2 M_{\mathrm{drop},i}\,e^{-s_\sigma^2/2}
\]

\subsubsection*{Step 7. Assembly}

The common base density $\phi_0$ appears in both $\widehat\phi_i$ and $\phi^*$ and cancels. Adding Steps 4, 5 and 6 by the triangle inequality and restricting the global $L^1$ estimate to $B(p_i,R_\ell)\cap D$,

\begin{align*}
  \bigl\|\widehat\phi_i-\phi^*\bigr\|_{L^1} &\le \bigl\|K_\sigma*\widehat\nu_i - K_\sigma*\nu_h^*\bigr\|_1 + \bigl\|K_\sigma*\nu_h^* - K_\sigma*\nu^*\bigr\|_1 + \varepsilon_{\mathrm{tail},i} && \text{\scriptsize triangle inequality} \\
   &\le L_{K,1}P_{\max}\varepsilon_{x,i} + 2\pi\sigma^2 E_{w,i} + \tfrac12 L_{K,1}L_T P_{\max}h_s + \varepsilon_{\mathrm{tail},i} && \text{\scriptsize Steps 4-6} \\
   &= L_{K,1}P_{\max}\!\left(\varepsilon_{\xi,i}+\sqrt2\,v + \frac{L_T h_s}{2}\right) + 2\pi\sigma^2 E_{w,i} + \varepsilon_{\mathrm{tail},i} \;=\; \varepsilon_{\phi,i} && \text{\scriptsize definition of $\varepsilon_{x,i}$}
\end{align*}

which is (L1). $\qquad\blacksquare$

\subsection{Evidence boundary}

Wolfram checked the two kernel integrals exactly, and the sup-norm, the normal-angle identity and the bounded-step inequality symbolically. The pushforward construction, the translation inequality (T), the curve quadrature, the atom matching of Step 5 and the tail decomposition are the analytic argument above; no computer-algebra output is offered in their place.

\newpage

\section{Lemma 2 - mass, quadrature and centroid perturbation}

\subsection{Statement}

Assume A1-A6 and the exact-cell contract $R_{\mathrm{comm}} \ge 2R_\ell$. Put

\[
r_0 = \min\Bigl\{R_\ell,\ \tfrac{d_s}{2},\ \tfrac{W_D}{2},\ \tfrac{H_D}{2}\Bigr\},\qquad m_{\min} = \frac{\phi_0\pi r_0^2}{4} > 0,\qquad m_{\max} = \pi R_\ell^2\bigl(\phi_0+P_{\max}\bigr)
\]

\[
\widehat\phi_{\max} = \phi_0+P_{\max}+E_{w,i},\qquad \widehat L_\phi \le \bigl(P_{\max}+E_{w,i}\bigr)\frac{e^{-1/2}}{\sigma}
\]

For a midpoint mesh of side at most $h$ set $r_h = h/\sqrt2$ (the half-diagonal), $A_V \le \pi R_\ell^2$ and $P_V \le 2\pi R_\ell$, and define

\[
\eta_m = \widehat L_\phi\, r_h A_V + \widehat\phi_{\max}\bigl(2P_V r_h + \pi r_h^2\bigr)
\]

\[
\eta_a = \bigl(\widehat\phi_{\max}+R_\ell \widehat L_\phi\bigr) r_h A_V + R_\ell\widehat\phi_{\max}\bigl(2P_V r_h + \pi r_h^2\bigr)
\]

If the gate $\eta_m < m_{\min}$ holds, set

\[
\varepsilon_q = \frac{\eta_a + R_\ell\,\eta_m}{m_{\min}-\eta_m}
\]

Then the implementation centroid obeys

\begin{equation*}
\bigl\|\widehat c_i - c_i^*\bigr\| \ \le\ \varepsilon_{c,i} \ :=\ \frac{2R_\ell\,\varepsilon_{\phi,i}}{m_{\min}} + \varepsilon_{q,i} + \varepsilon_{G,i}
\tag{L2}
\end{equation*}

In strict synchronised theorem mode $\varepsilon_{G,i}=0$. Delayed or missing neighbour positions require a separate geometric cell-error certificate; that error must not be hidden inside $\varepsilon_{\phi,i}$.

\subsection{Proof}

\subsubsection*{Step 1. Mass floor}

First, separation alone puts a whole disk inside agent $i$'s Voronoi cell. Let $\|q-p_i\| < d_s/2$. For any $j \ne i$,

\[
\|q-p_j\| \ \ge\ \|p_i-p_j\| - \|q-p_i\| \ \ge\ d_s - \frac{d_s}{2} \ =\ \frac{d_s}{2} \ >\ \|q-p_i\|
\]

so $q$ is strictly closer to $p_i$, i.e. $B(p_i,d_s/2)\subset \mathrm{Vor}_i(P)$. Intersecting with the local disk and the rectangle, $V_i^\ell \supseteq B(p_i,r_0)\cap D$. The worst placement of $p_i$ in a rectangle is a corner, which retains a quarter disk, so $|V_i^\ell| \ge \pi r_0^2/4$. Since $\phi^* \ge \phi_0$,

\[
m_i^* = \int_{V_i^\ell}\phi^*\,dq \ \ge\ \phi_0\cdot\frac{\pi r_0^2}{4} \ =\ m_{\min} \ >\ 0
\]

\begin{dbcallout}{NON-RECTANGULAR WORKSPACES}
The quarter-disk constant is specific to a rectangle. For another workspace, replace it with a uniform interior-cone constant; the quarter-disk value must not be carried over.
\end{dbcallout}

\subsubsection*{Step 2. Mass ceiling and mapped-density bounds}

Since $K_\sigma \le 1$ pointwise and $P_\Gamma \le P_{\max}$,

\[
\phi^*(q,t) = \phi_0 + \int_{\Gamma(t)}K_\sigma(q-\xi)\,ds \ \le\ \phi_0 + P_\Gamma \ \le\ \phi_0 + P_{\max}
\]

and $|V_i^\ell| \le |B(p_i,R_\ell)| = \pi R_\ell^2$, hence

\[
m_i^* \ \le\ \pi R_\ell^2\bigl(\phi_0+P_{\max}\bigr) \ =\ m_{\max}
\]

The measure decomposition of Lemma 1 bounds the mapped total weight by $P_{\max}+E_w$, which gives $\widehat\phi_{\max}$ by the same argument. Differentiating under the integral sign and using $\|\nabla K_\sigma\|_\infty = e^{-1/2}/\sigma$ gives the Lipschitz constant $\widehat L_\phi$.

\subsubsection*{Step 3. Midpoint mass error}

Split the midpoint cells into those lying wholly inside $V_i^\ell$ and those cut by its boundary. For an interior cell $Q$ with midpoint $q_Q$, Lipschitz continuity gives

\[
\left|\int_Q \phi\,dq - \phi(q_Q)\,|Q|\right| \le \widehat L_\phi \max_{q\in Q}\|q-q_Q\|\,|Q| = \widehat L_\phi\, r_h\,|Q|
\]

because the farthest point of a cell from its midpoint is a corner, at distance $r_h = h/\sqrt2$. Summing over interior cells contributes at most $\widehat L_\phi r_h A_V$.

The cut cells all lie within distance $r_h$ of $\partial V_i^\ell$. $V_i^\ell$ is convex, being an intersection of half-planes, a disk and a rectangle, so the inner and outer parallel sets of that width have combined area at most $2P_V r_h + \pi r_h^2$. Bounding the density there by $\widehat\phi_{\max}$ and adding the two groups gives $\eta_m$. Convexity also justifies $P_V \le 2\pi R_\ell$, the perimeter of the enclosing disk.

\subsubsection*{Step 4. First-moment error}

Repeat Step 3 for the \emph{shifted} first moment, which is what the centroid actually needs:

\[
a = \int_{V_i^\ell}(q-p_i)\,\phi(q)\,dq
\]

Its integrand $f(q) = (q-p_i)\phi(q)$ satisfies, on the cell,

\[
\|f(q)\| \le R_\ell\,\widehat\phi_{\max},\qquad \mathrm{Lip}(f) \le \widehat\phi_{\max} + R_\ell \widehat L_\phi
\]

since $\|q-p_i\| \le R_\ell$ there. Substituting these two constants for $\widehat L_\phi$ and $\widehat\phi_{\max}$ in Step 3 yields $\|\widehat a - a\| \le \eta_a$.

\subsubsection*{Step 5. Ratio perturbation}

The centroid is a ratio, so the two errors must be combined through a quotient bound rather than added. Write $\widehat c - p_i = \widehat a/\widehat m$ and $c - p_i = a/m$. Then

\[
\frac{\widehat a}{\widehat m} - \frac{a}{m} = \frac{\widehat a - a}{\widehat m} + a\left(\frac{1}{\widehat m}-\frac{1}{m}\right) = \frac{\widehat a - a}{\widehat m} + \frac{a\,(m-\widehat m)}{m\,\widehat m}
\]

The gate $\eta_m < m_{\min}$ makes the denominator safely positive, $\widehat m \ge m - \eta_m \ge m_{\min}-\eta_m > 0$, and $\|a\| \le R_\ell m$ because $\|q-p_i\|\le R_\ell$ on the cell. Therefore

\begin{align*}
  \left\|\frac{\widehat a}{\widehat m}-\frac{a}{m}\right\| &\le \frac{\|\widehat a-a\|}{\widehat m} + \frac{\|a\|}{m}\cdot\frac{|m-\widehat m|}{\widehat m} && \text{\scriptsize triangle inequality} \\
   &\le \frac{\eta_a}{m_{\min}-\eta_m} + R_\ell\,\frac{\eta_m}{m_{\min}-\eta_m} && \text{\scriptsize Steps 3-4 and $\|a\|\le R_\ell m$} \\
   &= \frac{\eta_a + R_\ell \eta_m}{m_{\min}-\eta_m} = \varepsilon_q && \text{\scriptsize definition}
\end{align*}

\subsubsection*{Step 6. Density perturbation on the same cell}

Now hold the cell fixed and change only the density. Lemma 1 gives

\[
\bigl|\widetilde m - m^*\bigr| = \left|\int_{V_i^\ell}(\widehat\phi_i-\phi^*)\,dq\right| \le \varepsilon_\phi,\qquad \bigl\|\widetilde a - a^*\bigr\| = \left\|\int_{V_i^\ell}(q-p_i)(\widehat\phi_i-\phi^*)\,dq\right\| \le R_\ell\,\varepsilon_\phi
\]

Both densities keep the base-density floor, so both masses are at least $m_{\min}$. Applying the same ratio identity as in Step 5,

\[
\bigl\|\widetilde c - c^*\bigr\| \le \frac{R_\ell\varepsilon_\phi}{m_{\min}} + R_\ell\cdot\frac{\varepsilon_\phi}{m_{\min}} = \frac{2R_\ell\,\varepsilon_\phi}{m_{\min}}
\]

\subsubsection*{Step 7. Assembly}

Adding the density error of Step 6, the midpoint error of Step 5 and the separately certified geometric cell error gives (L2). $\qquad\blacksquare$

The denominator gate is fail-closed in the certificate generator: a configuration with $\eta_m \ge m_{\min}$ produces no certificate at all rather than a large one.

\newpage

\section{Proposition 1 - ideal Local-CVT gradient and descent}

\subsection{Statement}

At every regular configuration covered by A1-A4 and A9,

\begin{equation*}
\nabla_{p_i}H^*(P,t) = 2\,m_i^*\,\bigl(p_i - c_i^*\bigr)
\tag{P1-a}
\end{equation*}

so the ideal Lloyd command is a scaled gradient step,

\[
u_i^* = -k_c\bigl(p_i-c_i^*\bigr) = -\frac{k_c}{2m_i^*}\,\nabla_{p_i}H^*
\]

and, with $A(P,t) = \mathrm{diag}_i\bigl(\tfrac{k_c}{2m_i^*}I_2\bigr)$,

\begin{equation*}
A(P,t) \succeq a_0 I,\quad a_0 = \frac{k_c}{2m_{\max}},\qquad \dot H^* = -\nabla H^{*\top}A\,\nabla H^* \le -a_0\bigl\|\nabla H^*\bigr\|^2
\tag{P1-b}
\end{equation*}

the last identity holding with the reference frozen.

\subsection{Proof}

\subsubsection*{Step 1. Partition the integral}

Away from the zero-measure tie and threshold sets of A9, the domain splits into the local cells and the saturated remainder $U = D\setminus\bigcup_j B(p_j,R_\ell)$, on which the running minimum equals $R_\ell^2$:

\[
H^* = \sum_j \int_{V_j^\ell}\|q-p_j\|^2\,\phi^*\,dq \;+\; R_\ell^2\int_U \phi^*\,dq
\]

\subsubsection*{Step 2. Volume term}

Differentiating the integrand of the $j=i$ term with respect to $p_i$,

\[
\frac{\partial}{\partial p_i}\|q-p_i\|^2 = 2\,(p_i-q) \quad\Longrightarrow\quad \int_{V_i^\ell}2(p_i-q)\phi^*\,dq = 2\left[m_i^* p_i - \int_{V_i^\ell}q\,\phi^*dq\right] = 2m_i^*\bigl(p_i-c_i^*\bigr)
\]

No other volume term depends on $p_i$: for $j\ne i$ the integrand $\|q-p_j\|^2$ does not involve $p_i$, and the $U$ integrand is the constant $R_\ell^2$.

\subsubsection*{Step 3. The moving-boundary terms cancel}

This is the step that the saturated cost exists to make work. Moving $p_i$ deforms $\partial V_i^\ell$, so Reynolds' transport theorem contributes a flux for each moving piece of boundary. There are three kinds, and each vanishes for its own reason.

\begin{itemize}[leftmargin=1.2em,itemsep=2pt]
  \item \emph{Shared Voronoi faces.} On the face between cells $i$ and $j$ we have $\|q-p_i\| = \|q-p_j\|$, so the two cells carry the same integrand value there. Whatever area cell $i$ gains, cell $j$ loses, with equal and opposite normal velocity; the two fluxes cancel exactly.
  \item \emph{The moving circle} $\|q-p_i\| = R_\ell$. Just inside, the integrand is $\|q-p_i\|^2 = R_\ell^2$; just outside, the saturated value is also $R_\ell^2$. The integrand is continuous across the circle, so the flux cancels. Without saturation the outer value would be $0$, the jump would be $R_\ell^2$, and a term proportional to $R_\ell^2\oint\phi^*$ would survive.
  \item \emph{The workspace boundary} $\partial D$ is fixed, so its normal velocity is zero.
\end{itemize}

Only the volume derivative survives, which is (P1-a).

\subsubsection*{Step 4. Descent}

By Lemma 2, $m_i^* \le m_{\max}$, so every diagonal block satisfies $\tfrac{k_c}{2m_i^*} \ge \tfrac{k_c}{2m_{\max}} = a_0$, i.e. $A \succeq a_0 I$. Substituting $\dot P = -A\nabla H^*$ with the reference frozen,

\[
\dot H^* = \nabla H^{*\top}\dot P = -\nabla H^{*\top}A\nabla H^* \le -a_0\|\nabla H^*\|^2 \ \le\ 0
\]

\subsubsection*{Step 5. The implemented cell is exact, not approximate}

Suppose $q \in B(p_i,R_\ell)$ and agent $j$ beats agent $i$ at $q$, i.e. $\|q-p_j\| \le \|q-p_i\| \le R_\ell$. Then

\[
\|p_j-p_i\| \le \|p_j-q\| + \|q-p_i\| \le 2R_\ell \ \le\ R_{\mathrm{comm}}
\]

so every possible local competitor is already inside the communication neighbourhood. Under A2 the implemented cell therefore \emph{equals} $V_i^\ell$; this is an identity, not an approximation. $\qquad\blacksquare$

\begin{dbcallout}{IF $R_{\mathrm{comm}} = R_\ell$}
The argument above needs $2R_\ell$. A theorem-mode configuration with $R_{\mathrm{comm}} < 2R_\ell$ is rejected by the validator, and the exact-cell claim would have to be marked BLOCKED rather than weakened silently.
\end{dbcallout}

Wolfram checked the integrand derivative and one exact fixed-square example. Those are sanity checks; the shape-derivative cancellation in Step 3 is the proof.

\newpage

\section{Proposition 2 - aggregate implementation disturbance}

\subsection{Statement}

Write the realised stacked dynamics as an ideal field plus a disturbance:

\[
\dot P = -A(P,t)\,\nabla H^*(P,t) + d(t),\qquad \|d(t)\|\le\delta
\]

\begin{equation*}
\delta = k_c\sqrt{\sum_{i=1}^N \varepsilon_{c,i}^2} + \varepsilon_{\mathrm{zoh}} + \varepsilon_{\mathrm{num}} + \varepsilon_{\mathrm{trk}} + \varepsilon_{\mathrm{sf}}
\tag{P2}
\end{equation*}

\subsection{Proof}

\subsubsection*{Step 1. Centroid component}

The nominal law and the ideal field differ only through the centroid each uses, so the gain multiplies the Lemma 2 error directly:

\[
u_i^{\mathrm{nom}} - u_i^* = -k_c\bigl(p_i-\widehat c_i\bigr) + k_c\bigl(p_i-c_i^*\bigr) = k_c\bigl(\widehat c_i - c_i^*\bigr) \quad\Longrightarrow\quad \bigl\|u_i^{\mathrm{nom}}-u_i^*\bigr\| \le k_c\,\varepsilon_{c,i}
\]

Stacking over the $N$ agents in the frozen product norm,

\[
\Bigl\|\bigl(k_c(\widehat c_i-c_i^*)\bigr)_{i=1}^N\Bigr\| = k_c\sqrt{\sum_i\bigl\|\widehat c_i-c_i^*\bigr\|^2} \le k_c\sqrt{\sum_i \varepsilon_{c,i}^2}
\]

\subsubsection*{Step 2. Remaining deviations}

Sample-and-hold, solver/integration, low-level tracking and safety-projection deviations are each bounded in the same stacked norm and added by the triangle inequality, which gives (P2).

\subsubsection*{Step 3. Analytic ceiling for the safety filter}

The safety term admits a bound that needs no model of the filter at all, only the speed cap that both the nominal and the filtered command respect. If $\|u_i\| \le u_{\max}$ for both, then

\[
\bigl\|u_i^{\mathrm{sf}}-u_i^{\mathrm{nom}}\bigr\| \le \bigl\|u_i^{\mathrm{sf}}\bigr\| + \bigl\|u_i^{\mathrm{nom}}\bigr\| \le 2u_{\max} \quad\Longrightarrow\quad \varepsilon_{\mathrm{sf}} \le \sqrt{N\,(2u_{\max})^2} = 2\sqrt{N}\,u_{\max}
\]

The norm convention is read from the implementation, never guessed from the symbol. If the cap were componentwise, $\|u_i\|_\infty \le u_{\max}$, then $\|u_i\|_2 \le \sqrt2 u_{\max}$ and the stacked ceiling would instead be $2\sqrt{2N}\,u_{\max}$. $\qquad\blacksquare$

\begin{dbcallout}{THIS CEILING IS DELIBERATELY CRUDE}
$2\sqrt N u_{\max}$ assumes the filter may reverse every command. A measured tightening may replace it only when its provenance and coverage interval are recorded in the ledger. Any theorem-mode gap, frontier or lead-side bias must be zero or appear as a further term in (P2).
\end{dbcallout}

\newpage

\section{Proposition 3 - reference rate and local regularity}

\subsection{Target and density rate}

Differentiating the offset target $\xi = b + d\,n$ by the product rule,

\[
\dot\xi = \dot b + \dot d\,n + d\,\dot n
\]

Under the rigid-body certificate A3, a point of the boundary at radius at most $R_O$ from the object centre moves at speed at most $\bar v_O + \bar\omega_O R_O$, the unit normal rotates at rate at most $\bar\omega_O$, and the offset varies at rate at most $\bar\omega_O L_d$. Hence

\[
\|\dot\xi\| \le \bar v_O + \bar\omega_O\bigl(R_O + d_{\max} + L_d\bigr) =: v_{\xi,\max}
\]

with $L_d = 0$ for the constant offset used in version 1. Differentiating the density under the integral sign and applying (T),

\begin{equation*}
\partial_t\phi^*(q,t) = -\int_{\Gamma(t)}\nabla K_\sigma\bigl(q-\xi(s,t)\bigr)\cdot\dot\xi(s,t)\,ds \quad\Longrightarrow\quad \nu_\phi := \bigl\|\partial_t\phi^*\bigr\|_1 \le L_{K,1}P_{\max}\,v_{\xi,\max}
\tag{P3-a}
\end{equation*}

\subsection{Cost and optimal-value rate}

Because the saturated integrand is bounded, $0 \le \min_i F_R \le R_\ell^2$, the cost rate inherits the density rate:

\[
\bigl|\partial_t H^*(P,t)\bigr| = \left|\int_D \min_i F_R\bigl(\|q-p_i\|\bigr)\,\partial_t\phi^*(q,t)\,dq\right| \le R_\ell^2\int_D\bigl|\partial_t\phi^*\bigr|dq = R_\ell^2\,\nu_\phi
\]

The optimal value inherits it too. For a fixed compact $X$ and any two times, the elementary inequality $\bigl|\min f - \min g\bigr| \le \sup|f-g|$ gives

\[
\bigl|H^*_{\min}(t_2)-H^*_{\min}(t_1)\bigr| \le \sup_{P\in X}\bigl|H^*(P,t_2)-H^*(P,t_1)\bigr| \le R_\ell^2\nu_\phi\,|t_2-t_1|
\]

so $H^*_{\min}$ is Lipschitz and, being Lipschitz, differentiable almost everywhere with $\bigl|\dot H^*_{\min}\bigr| \le R_\ell^2\nu_\phi$. Adding the two contributions bounds the drift of $V = H^*-H^*_{\min}$:

\begin{equation*}
\bigl|\partial_t H^* - \dot H^*_{\min}\bigr| \le \nu_H := 2R_\ell^2\,\nu_\phi
\tag{P3-b}
\end{equation*}

\begin{dbcallout}{RIGID BOUNDARIES ONLY}
This derivation uses a rigid twist. Deformation, a time-varying perimeter, or split/merge events break it; those cases need a direct measure-rate assumption on $\partial_t\nu_t^*$ instead.
\end{dbcallout}

\subsection{Local PL and quadratic growth (M5)}

Suppose A8 holds: on a convex, symmetry-reduced chart of the selected branch, $\nabla^2 H^*_{\mathrm{red}}(z,t) \succeq m_H I$ with $m_H>0$. Write $f$ for $H^*$ in those coordinates and $z^*$ for a minimiser. Integrating the Hessian twice along the segment from $z^*$ to $z$ gives the standard strong-convexity inequality

\[
f(y) \ \ge\ f(z) + \nabla f(z)^\top(y-z) + \frac{m_H}{2}\|y-z\|^2 \qquad \forall\,y,z
\]

Putting $y = z^*$ gives quadratic growth,

\begin{equation*}
H^* - H^*_{\min} \ \ge\ \frac{\alpha}{2}\,\operatorname{dist}^2\bigl(P,\mathcal{C}^*(t)\bigr),\qquad \alpha = m_H
\tag{QG}
\end{equation*}

and minimising the right-hand side over $y$ instead, at $y = z - \nabla f(z)/m_H$, gives

\begin{align*}
  H^*_{\min} \ &\ge\ H^*(z) - \frac{\|\nabla f(z)\|^2}{2m_H} && \text{\scriptsize minimise the quadratic lower bound over $y$} \\
  \tfrac12\bigl\|\nabla H^*\bigr\|^2 \ &\ge\ \mu\bigl[H^*-H^*_{\min}\bigr],\qquad \mu = m_H && \text{\scriptsize rearrange}
\end{align*}

which is the local Polyak-Lojasiewicz inequality used in Theorem 1. $\qquad\blacksquare$

\begin{dbcallout}{A8 IS CONDITIONAL, AND THAT IS THE HONEST STATUS}
The implication above is proved. Its premise is not established here: no specific nondegenerate branch, cell combinatorics, reduced gauge, or complete parameter tube was supplied, so no validated interval enclosure of the reduced Hessian exists. A floating-point eigenvalue at one configuration would be NUMERIC-SANITY evidence and is not presented as proof. No global PL claim is made for a nonconvex CVT cost.
\end{dbcallout}

\newpage

\section{Theorem 1 - local practical stability}

\subsection{Statement}

Assume A1-A10, Lemmas 1-2 and Propositions 1-3, on an interval throughout which the solution remains in the regular tube $X$. Define

\[
V(P,t) = H^*(P,t)-H^*_{\min}(t)\ \ge\ 0,\qquad a_0 = \frac{k_c}{2m_{\max}},\qquad \lambda_H = a_0\mu,\qquad B = \frac{\delta^2}{2a_0} + \nu_H
\]

Then for every $t \ge t_0$ in that interval

\begin{equation*}
V(t)\ \le\ e^{-\lambda_H(t-t_0)}V(t_0) + \frac{B}{\lambda_H}\Bigl(1-e^{-\lambda_H(t-t_0)}\Bigr)
\tag{T1-a}
\end{equation*}

and, if the regular-tube hypotheses persist for all future time,

\begin{equation*}
\limsup_{t\to\infty}\operatorname{dist}\bigl(P(t),\mathcal{C}^*(t)\bigr) \ \le\ \sqrt{\frac{2B}{\alpha\,\lambda_H}}
\tag{T1-b}
\end{equation*}

\subsection{Proof - every inequality named}

Differentiate $V$ along the solution at an almost-every regular time and bound each term by the result that produces it.

\begin{align*}
  \dot V &= \nabla H^{*\top}\dot P + \partial_t H^* - \dot H^*_{\min} && \text{\scriptsize chain rule; $H^*_{\min}$ differentiable a.e. by Prop. 3} \\
   &= \nabla H^{*\top}\bigl(-A\nabla H^* + d\bigr) + \partial_t H^* - \dot H^*_{\min} && \text{\scriptsize substitute the dynamics} \\
   &\le -\nabla H^{*\top}A\,\nabla H^* + \bigl\|\nabla H^*\bigr\|\,\|d\| + \nu_H && \text{\scriptsize Cauchy-Schwarz, and (P3-b)} \\
   &\le -a_0\bigl\|\nabla H^*\bigr\|^2 + \delta\bigl\|\nabla H^*\bigr\| + \nu_H && \text{\scriptsize $A \succeq a_0 I$ by (P1-b); $\|d\|\le\delta$ by (P2)} \\
   &\le -a_0\bigl\|\nabla H^*\bigr\|^2 + \frac{a_0}{2}\bigl\|\nabla H^*\bigr\|^2 + \frac{\delta^2}{2a_0} + \nu_H && \text{\scriptsize Young: $\delta x \le \tfrac{a_0}{2}x^2 + \tfrac{\delta^2}{2a_0}$} \\
   &= -\frac{a_0}{2}\bigl\|\nabla H^*\bigr\|^2 + B && \text{\scriptsize definition of $B$} \\
   &\le -a_0\mu\,V + B \;=\; -\lambda_H V + B && \text{\scriptsize local PL: $\tfrac12\|\nabla H^*\|^2 \ge \mu V$}
\end{align*}

The scalar comparison lemma applied to $\dot V \le -\lambda_H V + B$ integrates this to (T1-a): the solution of the corresponding equality is exactly the right-hand side of (T1-a), and a differential inequality is dominated by it. Letting $t\to\infty$ in (T1-a), the exponential terms vanish and

\[
\limsup_{t\to\infty} V(t) \ \le\ \frac{B}{\lambda_H}
\]

Finally, quadratic growth (QG) converts the Lyapunov level into a distance,

\[
\operatorname{dist}\bigl(P,\mathcal{C}^*\bigr) \le \sqrt{\frac{2V}{\alpha}} \quad\Longrightarrow\quad \limsup_{t\to\infty}\operatorname{dist}\bigl(P(t),\mathcal{C}^*(t)\bigr) \le \sqrt{\frac{2B}{\alpha\lambda_H}}
\]

which is (T1-b). $\qquad\blacksquare$

\subsection{Status}

The implication is analytically complete and its scalar algebra is Wolfram-verified (Young, the comparison solution, and the strong convexity to PL step). The theorem remains CONDITIONAL on A5 and A8-A10 and on the run/tube certificates. It is local. It is not a global PL theorem, and it says nothing about caging or transport success.

\newpage

\section{Corollaries}

\subsection{Frozen perfect information}

If every implementation error vanishes and the reference is frozen, then $\delta = 0$ and $\nu_H = 0$, so $B = 0$ and (T1-a) collapses to pure exponential decay:

\[
\dot V \le -\lambda_H V \quad\Longrightarrow\quad V(t) \le e^{-\lambda_H(t-t_0)}V(t_0) \quad\Longrightarrow\quad \operatorname{dist}\bigl(P(t),\mathcal{C}^*\bigr)\to 0
\]

This is a sanity gate, not a decoration: the constant generator and the unit tests require exactly $B = 0$ and a zero ultimate radius in this case. If they did not, the formulas or their wiring would be wrong and Theorem 1 could not be marked complete.

\subsection{Sampled-data corollary (M7)}

Let the implementation take explicit steps of size $\Delta t$,

\[
P_{k+1} = P_k + \Delta t\bigl[-A_k g_k + d_k\bigr],\qquad g_k = \nabla H^*(P_k,t_k),\qquad a_0 I \preceq A_k \preceq a_1 I,\quad a_1 = \frac{k_c}{2m_{\min}}
\]

with $\|d_k\| \le \delta$ and $\nabla H^*$ spatially $L_H$-Lipschitz on $X$. Define

\[
\bar a = a_0 - \frac{L_H a_1^2 \Delta t}{2},\qquad C_\delta = \delta\bigl(1 + L_H a_1 \Delta t\bigr),\qquad B_d = \frac{C_\delta^2}{2\bar a} + \frac{L_H}{2}\delta^2\Delta t + \nu_H
\]

Then, under the step gates

\begin{equation*}
\Delta t < \frac{2a_0}{L_H a_1^2},\qquad 0 < \bar a\,\mu\,\Delta t \le 1
\tag{M7-gate}
\end{equation*}

the Lyapunov level contracts geometrically,

\begin{equation*}
V_{k+1} \le \bigl(1-\bar a\mu\Delta t\bigr)V_k + B_d\,\Delta t,\qquad \limsup_{k\to\infty}\operatorname{dist}\bigl(P_k,\mathcal{C}_k^*\bigr) \le \sqrt{\frac{2B_d}{\alpha\,\bar a\,\mu}}
\tag{M7}
\end{equation*}

\subsubsection*{Proof}

Apply the descent lemma to the $L_H$-Lipschitz gradient, add the reference drift over one step, and write $x = \|g_k\|$:

\begin{align*}
  V_{k+1}-V_k &\le g_k^\top\bigl(P_{k+1}-P_k\bigr) + \frac{L_H}{2}\bigl\|P_{k+1}-P_k\bigr\|^2 + \nu_H\Delta t && \text{\scriptsize descent lemma and (P3-b)} \\
   &\le -a_0\Delta t\,x^2 + \delta\Delta t\,x + \frac{L_H}{2}\Delta t^2\bigl(a_1 x + \delta\bigr)^2 + \nu_H\Delta t && \text{\scriptsize $a_0 I \preceq A_k \preceq a_1 I$, $\|d_k\|\le\delta$} \\
   &= \Delta t\left[-\bar a\,x^2 + C_\delta\,x + \frac{L_H}{2}\delta^2\Delta t + \nu_H\right] && \text{\scriptsize expand the square and collect powers of $x$} \\
   &\le \Delta t\left[-\frac{\bar a}{2}x^2 + \frac{C_\delta^2}{2\bar a} + \frac{L_H}{2}\delta^2\Delta t + \nu_H\right] && \text{\scriptsize Young, with $\bar a > 0$} \\
   &= \Delta t\left[-\frac{\bar a}{2}x^2 + B_d\right] && \text{\scriptsize definition of $B_d$} \\
   &\le \Delta t\bigl[-\bar a\mu V_k + B_d\bigr] && \text{\scriptsize local PL: $\tfrac12 x^2 \ge \mu V_k$}
\end{align*}

Rearranging gives the recursion in (M7). The first gate is exactly the condition $\bar a > 0$ needed for the Young step; the second keeps the contraction factor in $[0,1)$. Summing the resulting geometric series,

\[
\limsup_{k\to\infty}V_k \le \frac{B_d\Delta t}{\bar a\mu\Delta t} = \frac{B_d}{\bar a\mu}
\]

and quadratic growth converts this to the stated ultimate radius. $\qquad\blacksquare$

Wolfram expanded the descent algebra and checked the coefficient collection; a randomised positive-parameter test independently compares the certified bound against the direct one-step expression.

\newpage

\section{Limits and evidence status}

\begin{center}
\small
\begin{longtable}{>{\RaggedRight\arraybackslash}p{3.1cm} >{\RaggedRight\arraybackslash}p{2.6cm} >{\RaggedRight\arraybackslash}p{8.2cm}}
\toprule
\textbf{Result} & \textbf{Status} & \textbf{Limiting condition} \\
\midrule
\endhead
Lemma 1 & PROVED + CAS & A5 map bounds and tail separation \\
Lemma 2 & PROVED + CAS & rectangle, midpoint gate, exact-cell contract \\
Proposition 1 & PROVED & regular tie/threshold sets; saturated cost \\
Proposition 2 & PROVED & declared stacked disturbance components \\
Proposition 3 & CONDITIONAL & rate proved; A8 branch coercivity assumed \\
Theorem 1 & CONDITIONAL & complete implication under A1-A10 \\
Sampled corollary & PROVED & explicit step gates \\
\bottomrule
\end{longtable}
\end{center}

\subsection{Evidence classes}

\begin{itemize}[leftmargin=1.2em,itemsep=2pt]
  \item PROVED: the analytic implication is written out in this supplement.
  \item VERIFIED: Wolfram reproduced the identified integral or algebraic subclaim, and nothing beyond it.
  \item CONDITIONAL: a named branch, tube, map bound, or runtime deviation certificate is still required.
  \item NUMERIC-SANITY: debugging evidence only. No such result is used as proof anywhere in this package, and there are no such entries.
\end{itemize}

\subsection{Fail-closed runtime gates}

Each of these refuses to emit a certificate rather than emitting a degraded one, and each is covered by a test that feeds it a deliberately invalid configuration:

\begin{itemize}[leftmargin=1.2em,itemsep=2pt]
  \item locality: a theorem-mode configuration with $R_{\mathrm{comm}} < 2R_\ell$ is rejected;
  \item arc length: a boundary observation without positive $\Delta s_k$ is rejected, and the unit-weight fallback is forbidden;
  \item quadrature rule: anything other than a midpoint grid is rejected;
  \item quadrature gate: $\eta_m \ge m_{\min}$ is rejected;
  \item bias: a nonzero gap or exploration gain is rejected unless separately budgeted into $\delta$;
  \item sampled step: $\Delta t$ violating (M7-gate) is rejected.
\end{itemize}

\subsection{Claim audit}

This supplement contains no global PL statement, no formal caging claim, no finite-time phase-completion result, and no unconditional transport or safety guarantee. The numeric certificate in Appendix A validates the dependency wiring and the gates; its ultimate radius is deliberately conservative and may be numerically vacuous. It is not experimental evidence and is not used to claim transport performance.

\newpage

\appendix

\section{Constant ledger}

Every value below is produced by one code path from the frozen formulas; none is entered by hand. The chain from raw sensing bounds to the ultimate radius is

\[
\bigl(\varepsilon_b,\varepsilon_n,\varepsilon_d,v,h_s,E_w,M_{\mathrm{drop}}\bigr) \ \longrightarrow\ \varepsilon_\phi \ \longrightarrow\ \varepsilon_c \ \longrightarrow\ \delta \ \longrightarrow\ B \ \longrightarrow\ \sqrt{\frac{2B}{\alpha\lambda_H}}
\]

\begin{center}
\small
\begin{longtable}{l r l l}
\toprule
\textbf{Constant} & \textbf{Value} & \textbf{Unit} & \textbf{Status} \\
\midrule
\endhead
kernel\_l1 & 0.251327 & m\textasciicircum{}2 & verified \\
kernel\_grad\_l1 & 1.57496 & m & verified \\
kernel\_grad\_linf & 3.03265 & 1/m & verified \\
epsilon\_xi & 0.00449999 & m & conditional \\
epsilon\_phi & 0.106034 & m\textasciicircum{}3 & proved \\
m\_min & 5.02655 & m\textasciicircum{}3 & proved \\
m\_max & 34.5827 & m\textasciicircum{}3 & proved \\
grid\_spacing & 0.000390625 & m & verified \\
eta\_m & 0.0598959 & m\textasciicircum{}3 & proved \\
eta\_a & 0.0574695 & m\textasciicircum{}4 & proved \\
epsilon\_q & 0.0212188 & m & proved \\
epsilon\_c & 0.0549703 & m & proved \\
delta & 0.156932 & m/s & conditional \\
nu\_phi & 0.464928 & m\textasciicircum{}3/s & proved \\
nu\_H & 0.595108 & m\textasciicircum{}5/s & proved \\
a0 & 0.0130123 & 1/(s*m\textasciicircum{}3) & proved \\
lambda\_H & 0.000650615 & 1/s & conditional \\
B & 1.54143 & m\textasciicircum{}5/s & proved \\
ultimate\_radius & 307.843 & m & conditional \\
sampled\_Bd & 1.54617 & m\textasciicircum{}5/s & proved \\
sampled\_ultimate\_radius & 308.792 & m & conditional \\
\bottomrule
\end{longtable}
\end{center}

\newpage

\section{Wolfram verification manifest}

Connector: Wolfram (version not exposed). Local runtime: WolframScript 1.14.0; all nine committed scripts executed with exit code 0. Exact inputs, assumptions and outputs are archived under theory/wolfram/results.

\begin{center}
\small
\begin{longtable}{>{\RaggedRight\arraybackslash}p{2.4cm} >{\RaggedRight\arraybackslash}p{6.2cm} >{\RaggedRight\arraybackslash}p{2.3cm} >{\RaggedRight\arraybackslash}p{2.6cm}}
\toprule
\textbf{ID} & \textbf{Target} & \textbf{Evidence} & \textbf{Verdict} \\
\midrule
\endhead
W-L1-01/02 & Gaussian $L^1$ integrals & EXACT & PASS \\
W-L1-03 & gradient radial maximum & SYMBOLIC & PASS \\
W-L1-04/05 & normal and bounded-step identities & SYMBOLIC & PASS \\
W-L2-01/02 & ratio, monotonicity, mass bounds & SYMBOLIC & PASS \\
W-P1-01 & integrand derivative, fixed-cell check & SYMBOLIC & PASS \\
W-P3-01 & reference-rate algebra & SYMBOLIC & PASS \\
W-T1-01 & Young and comparison ODE & SYMBOLIC & PASS \\
W-C1-01 & sampled-data descent expansion & SYMBOLIC & PASS \\
W-M5-01 & generic quadratic implication & SYMBOLIC & PASS \\
W-M5-02 & actual branch Hessian box & INTERVAL & CONDITIONAL \\
\bottomrule
\end{longtable}
\end{center}

Evidence rule: an exact or symbolic output supports only the stated integral or algebraic identity. Interval evidence would have to cover the complete selected tube to support A8, and none does.

\subsection{What computer algebra did not verify}

\begin{itemize}[leftmargin=1.2em,itemsep=2pt]
  \item the pushforward construction, the $L^1$ translation inequality, the curve quadrature and the atom matching of Lemma 1;
  \item the Voronoi shape derivative and the moving-boundary flux cancellation of Proposition 1;
  \item any uniform lower bound on a reduced Hessian for a named branch, which is why A8 and therefore Theorem 1 remain conditional.
\end{itemize}

\subsection{Reference}

J. Cortes, S. Martinez and F. Bullo, ``Spatially-distributed coverage optimization and control with limited-range interactions'', ESAIM: COCV, 11(4):691-719, 2005. DOI 10.1051/cocv:2005024.

\end{document}
